\chapter{Varicella-Zoster Virus (VZV)}

\section*{What Is This Test?}
Varicella-zoster virus (VZV) testing detects VZV, a member of the \textit{Herpesviridae} family (human herpesvirus 3/HHV-3). Primary VZV infection causes varicella (chickenpox), characterized by a generalized pruritic vesicular rash that appears in crops, starting on the trunk and face and spreading centrifugally, with lesions in various stages (papules, vesicles, pustules, crusts). After primary infection, VZV establishes latency in dorsal root ganglia. Reactivation later in life produces herpes zoster (shingles), a painful, unilateral, dermatomal vesicular eruption often preceded by prodromal pain, itching, or paresthesia in the affected dermatome \cite{gnann2002herpes}. Postherpetic neuralgia (PHN), defined as pain persisting >90 days after rash onset, is the most common complication, occurring in 10--18\% of zoster cases and increasing with age (up to 30\% in those over 80) \cite{cohen2013herpes}. Other complications include ocular involvement (herpes zoster ophthalmicus), Ramsay Hunt syndrome (facial nerve palsy with vesicles in the ear canal), bacterial superinfection, and, in immunocompromised patients, disseminated zoster with viremia and visceral involvement (pneumonitis, hepatitis, encephalitis).

\section*{Alternative Names}
Chickenpox Testing; Shingles Testing; VZV Serology; VZV PCR; VZV DFA; Varicella Immunity Panel; Varicella Titer; VZV IgG/IgM; Herpes Zoster Testing; Tzank Smear (historical)

\section*{Principle}
PCR: Detection of VZV DNA by real-time PCR from vesicular fluid swabs, lesion scrapings, CSF, or tissue. Sensitivity >95\% and specificity >99\%, making it the gold standard for acute diagnosis. Direct fluorescent antibody (DFA): Vesicle scraping smeared on a slide, fixed, and stained with fluorescein-labeled VZV-specific monoclonal antibodies; detects VZV-infected cells (multinucleated giant cells with fluorescent inclusions). Sensitivity 80--95\% compared to PCR; requires adequate cellularity. Results in 1--2 hours. Viral culture: Inoculation of vesicle fluid onto human fibroblast cell lines (MRC-5, WI-38). Cytopathic effect (focal rounding, multinucleated giant cells) appears in 3--7 days. Sensitivity 40--50\% compared to PCR; now rarely used for routine diagnosis. Serology: ELISA or CLIA detecting VZV IgG (past infection or vaccination) and VZV IgM (acute or recent infection). VZV IgG avidity testing can differentiate primary from past infection. Serology is primarily used for immunity screening, not for diagnosing active VZV infection (PCR and DFA are preferred for acute lesions). Tzank smear (Wright-Giemsa stain of vesicle scraping): Identifies multinucleated giant cells and ballooning degeneration of keratinocytes but cannot distinguish VZV from HSV; historically useful, now largely replaced by DFA and PCR.

\section*{History}
Varicella was distinguished from smallpox by William Heberden in 1767. The relationship between varicella and herpes zoster was recognized by James von Bokay in 1892. Thomas Weller isolated VZV in cell culture in 1953, proving that varicella and zoster are caused by the same virus. The live attenuated varicella vaccine (Varivax, Oka strain) was developed by Michiaki Takahashi in 1974 and licensed in the US in 1995. The varicella vaccine reduced incidence by 97\% in vaccinated populations \cite{marin2007prevention}. The herpes zoster vaccine (Zostavax, live attenuated) was licensed in 2006, and the recombinant zoster vaccine (Shingrix, adjuvanted VZV glycoprotein E) was licensed in 2017, demonstrating >90\% efficacy for prevention of herpes zoster and PHN in adults over 50 \cite{dooling2018recommendations}. PCR for VZV became available in the 1990s and has largely supplanted culture and DFA for diagnosis of CNS VZV disease (encephalitis, meningoencephalitis, vasculopathy) and atypical/ disseminated zoster in immunocompromised patients. The advent of antivirals (acyclovir 1980s; valacyclovir, famciclovir 1990s) provided effective treatment for both varicella and zoster when initiated early.

\section*{Why This Test Is Ordered}
\begin{itemize}
  \item Determination of varicella immunity status in healthcare workers, pregnant women, college students, military recruits, international travelers, and individuals prior to immunosuppressive therapy; documentation of positive VZV IgG is accepted as evidence of immunity per ACIP guidelines
  \item Prenatal screening for varicella immunity: identification of non-immune pregnant women who should receive postpartum vaccination (varicella vaccine is a live vaccine contraindicated during pregnancy); non-immune pregnant women exposed to varicella or zoster should receive VariZIG (varicella-zoster immune globulin) within 10 days of exposure
  \item Confirmation of suspected acute varicella (chickenpox) in atypical or breakthrough cases (vaccine-modified varicella, which may present with fewer than 50 lesions, mostly maculopapular rather than vesicular)
  \item Diagnosis of herpes zoster (shingles) when the presentation is atypical (zoster sine herpete, disseminated zoster, multidermatomal zoster in an immunocompromised host, or when the differential includes HSV)
  \item Diagnosis of VZV CNS disease: VZV PCR on CSF in patients with suspected VZV encephalitis, meningoencephalitis, vasculopathy, or myelitis, especially in immunocompromised hosts and the elderly \cite{gilden2013variegate}
  \item Evaluation of VZV retinitis (acute retinal necrosis or progressive outer retinal necrosis) by PCR of vitreous or aqueous fluid
  \item Confirmation of VZV pneumonitis in immunocompromised patients with pulmonary infiltrates and concomitant dermatomal or disseminated zoster
\end{itemize}

\section*{Primary vs Secondary Test}
VZV IgG is the primary test for immunity screening. For acute varicella or zoster diagnosis, PCR from vesicular lesions is the primary preferred diagnostic test due to superior sensitivity and specificity. DFA is a primary alternative when PCR is not available. Viral culture and Tzank smear are secondary/historical tests. For CNS disease, CSF VZV PCR is the primary diagnostic test. Serology (IgM) is a secondary test for acute infection, as VZV IgM has low sensitivity and specificity compared to PCR.

\section*{How This Test Is Performed}
Vesicle fluid/lesion specimen (PCR and DFA): Unroof a fresh vesicle with a sterile needle or scalpel blade, absorb the fluid with a flocked or Dacron swab, and vigorously swab the base of the lesion to collect infected epithelial cells. Place swab in viral transport medium (for PCR and culture) or roll/swipe directly onto a glass slide (for DFA). Crusted lesions have lower sensitivity. For DFA, the smear must contain adequate numbers of cells. Blood (serology): Venous blood collected in an SST tube; centrifuge and separate serum within 2 hours. CSF for VZV PCR: 0.5--1 mL in a sterile tube; transport to the laboratory immediately; do not refrigerate. Vitreous/aqueous fluid: collected by ophthalmologist.

\section*{What Preparation Is Needed}
No special patient preparation. For lesion specimens, collect before the application of topical antimicrobials, zinc oxide, or calamine lotion, as these may interfere with PCR. If a Tzank smear is included, the lesion should not be cleaned with alcohol before scraping, as this removes the cellular material needed for cytologic examination. Antiviral therapy (acyclovir, valacyclovir, famciclovir) started more than 24--48 hours before specimen collection may reduce PCR and DFA sensitivity.

\section*{Contraindications and Precautions}
VZV IgM may be falsely positive in the presence of rheumatoid factor or cross-reactivity with other herpesviruses (HSV, EBV, CMV). A positive VZV IgM in a patient with primary HSV-1 or HSV-2 infection may represent cross-reactivity rather than concurrent VZV infection; PCR of lesion specimens resolves this ambiguity. VZV IgG may become positive as early as 2--4 days after varicella rash onset, and a positive IgG in the setting of acute rash does NOT exclude primary VZV infection; detectable IgG at the time of rash onset may represent an anamnestic response from prior vaccination or remote infection. VZV PCR from skin lesions or CSF may be falsely negative if the specimen was inadequately collected (too few cells) or obtained >5 days after rash onset when viral load has declined. In cases of VZV CNS vasculopathy, PCR may be negative when disease is antibody-mediated rather than lytic, and detection of VZV IgG in CSF (with CSF:serum IgG ratio compared to albumin ratio) or anti-VZV IgM in CSF can aid diagnosis.

\section*{Risks}
Vesicle unroofing causes minor discomfort and may temporarily increase local pain. There is minimal risk of scarring at the unroofing site. Standard venipuncture risks apply for serology. CSF collection via lumbar puncture carries risks of post-LP headache, bleeding, or infection. No risks from the laboratory tests per se.

\section*{What Happens After the Test}
VZV IgG/IgM serology results available in 1--24 hours. PCR results available in 2--6 hours. DFA results available in 1--2 hours. Positive VZV IgM from a patient with suspected acute varicella should prompt placement in airborne and contact precautions (varicella is airborne) until all lesions are crusted over. Non-immune exposed individuals should be offered varicella vaccine (within 3--5 days of exposure for post-exposure prophylaxis) or VariZIG. For herpes zoster, the patient should keep lesions covered to prevent transmission of VZV to susceptible contacts. Antiviral therapy with valacyclovir (1 g PO TID x 7 days), famciclovir (500 mg PO TID x 7 days), or acyclovir (800 mg 5x daily x 7 days) should be initiated, ideally within 72 hours of rash onset. IV acyclovir (10 mg/kg q8h) is indicated for disseminated zoster, VZV encephalitis, or severe disease in immunocompromised hosts.

\section*{Understanding Results}

\subsection*{Below Normal}
\begin{itemize}
  \item \textbf{VZV IgG negative (nonreactive):} Indicates susceptibility to varicella---no evidence of past infection or vaccination. The individual should receive 2 doses of varicella vaccine (4--8 weeks apart) if no contraindications. This result is particularly important in pregnant women exposed to varicella, who should receive VariZIG and be counseled about the risks of congenital varicella syndrome
  \item \textbf{VZV PCR negative from lesion:} No VZV DNA detected. In a patient with a vesicular rash, this may indicate: HSV infection (HSV-1 or HSV-2 PCR should be performed); non-herpetic vesicular dermatosis; or a false-negative result from inadequate specimen collection (crust, insufficient cells) or late presentation (>5 days)
  \item \textbf{CSF VZV PCR negative:} Does not exclude VZV CNS disease. PCR sensitivity is lower in VZV vasculopathy compared to VZV encephalitis. CSF anti-VZV IgG index (Reiber diagram calculation) can aid in the diagnosis of VZV vasculopathy when PCR is negative
  \item \textbf{VZV IgM negative:} Does not exclude acute infection. IgM may be absent in breakthrough varicella, recurrent zoster, or CNS reactivation without skin lesions
\end{itemize}

\subsection*{Above Normal}
\begin{itemize}
  \item \textbf{VZV IgG positive (reactive):} Indicates immunity from past infection or vaccination. Generally protective against varicella. A positive IgG in a patient with a zosteriform rash is consistent with herpes zoster (reactivation in the setting of pre-existing immunity). Quantitative titers are not routinely followed; no specific cutoff correlates with protection against zoster
  \item \textbf{VZV IgM positive (reactive):} Indicates acute or recent VZV infection (primary varicella or recent reactivation). Not specific; false positives occur with other herpesvirus infections. Should be correlated with clinical presentation and, ideally, confirmed by PCR
  \item \textbf{VZV PCR positive from skin lesion:} Confirms active VZV infection. Cannot distinguish primary varicella from herpes zoster; this distinction is clinical (generalized versus dermatomal; age; prior history of varicella or vaccination). If the patient is immunocompromised, systemic (disseminated) disease should be considered
  \item \textbf{VZV PCR positive from CSF:} Confirms VZV CNS infection (encephalitis, meningitis, vasculopathy, myelitis). A positive PCR in the appropriate clinical context (fever, altered mental status, dermatomal zoster, stroke in a young patient) is diagnostic and should prompt IV acyclovir therapy
  \item \textbf{VZV PCR positive from vitreous/aqueous:} Confirms VZV retinitis (acute retinal necrosis or progressive outer retinal necrosis). Requires urgent ophthalmologic evaluation and combined systemic plus intravitreal antiviral therapy
\end{itemize}

\section*{Normal Results}
VZV IgG: \textbf{Positive (immune/reactive)} indicates past infection or vaccination and presumed protection against varicella. \textbf{Negative (nonreactive)} indicates susceptibility. VZV IgM: \textbf{Negative}. VZV PCR: \textbf{Not detected}.

\section*{Follow-up Tests}
\begin{itemize}
  \item \textbf{VZV PCR of vesicle fluid or lesion scraping:} Confirmatory test when DFA is negative but clinical suspicion remains or when atypical presentation (disseminated zoster, breakthrough varicella in a vaccinated patient) requires confirmation
  \item \textbf{HSV-1 and HSV-2 PCR of lesion specimen:} When the clinical presentation is vesicular but VZV PCR/DFA is negative, to rule out herpes simplex virus infection that may mimic zoster (especially sacral HSV mimicking sacral zoster, or orolabial HSV mimicking mandibular zoster)
  \item \textbf{CSF VZV PCR and CSF:serum VZV IgG ratio:} When VZV CNS disease is suspected but CSF PCR is negative; an elevated VZV IgG index (intrathecal production of anti-VZV antibodies) supports the diagnosis of VZV vasculopathy, myelitis, or encephalitis
  \item \textbf{CSF analysis (cell count, glucose, protein):} In suspected VZV CNS disease; may demonstrate lymphocytic pleocytosis (50--500 WBC/$\mu$L), normal to mildly reduced glucose, and elevated protein
  \item \textbf{Varicella vaccine:} For VZV IgG-negative individuals without contraindications; 2-dose series 4--8 weeks apart; post-vaccination serologic follow-up is not routinely recommended but may be considered for healthcare workers
  \item \textbf{Bacterial culture of lesions:} When zoster lesions appear superinfected (increased erythema, purulence, crusting with surrounding cellulitis) to guide antibiotic therapy for secondary bacterial infection (typically \textit{Staphylococcus aureus} or \textit{Streptococcus pyogenes})
\end{itemize}

\section*{References}
\bibliographystyle{unsrtnat}
\bibliography{references}

\clearpage
\section*{Regulatory Status and Authorization Date}
Regulatory status applies to the specific assay, instrument, manufacturer, and intended use rather than to the test name alone. No single approval date can be assigned to this test category. For a commercial assay, verify the current product in the relevant regulator's database: the U.S. Food and Drug Administration (FDA) clearance, approval, or authorization; India's Central Drugs Standard Control Organisation (CDSCO) licence or permission; the European Union In Vitro Diagnostic Regulation (EU IVDR) CE certificate issued through the applicable conformity-assessment route; or the United Kingdom Medicines and Healthcare products Regulatory Agency (MHRA) registration and UKCA/CE status. Laboratory-developed tests may not have an individual FDA, CDSCO, EU IVDR, or MHRA approval date.
\clearpage
